\section{Related Literature}
\paragraph{Empirical Evidence on Endorsements and Election Outcomes}
The role of elite endorsements in primary elections has been widely studied in political science and economics, with strong evidence that endorsements are a critical determinant of electoral outcomes. Early work from Cohen et al. \parencite*{cohen2003} established that endorsements are strongly correlated with nomination success, arguing that candidates with more elite backing secured more delegates at the nominating convention. Cohen et al. \parencite*{Cohen2008} expanded these results to argue that the party elite still maintain control of the nomination process, despite reforms in the 1970s meant to shift power to primary voters. Their finding highlights that endorsements are often as, if not more, effective at predicting nomination outcomes than polling numbers. 

Early models of presidential primary voting focused on factors such as candidate spending \parencite{Goldstein1978}, voter preferences (\cite{Marshall1983}; \cite{Norrander1986}) and expectations (\cite{Bartels1987}, \cite*{Bartels1988}), and polling trends (\cite{Mayer1996}, \cite*{Mayer2003}). These studies emphasized that candidate strength and media narratives influence voter decision-making. However, they largely overlooked the role of elite endorsements as a force shaping candidate viability. Steger and Davis \parencite*{Davis2000} address this gap by demonstrating that endorsement patterns significantly influence primary vote shares, reinforcing the idea that elite coordination serves as an informal mechanism of party control. Expanding on this work, Steger \parencite*{Steger2007} improves forecasting models of primary election outcomes by incorporating elite endorsements and perceptions of electability, showing that endorsement patterns influence voter behavior and nomination results.

However, the effectiveness of endorsements varies across elections. Steger \parencite*{Steger2008} finds that viability and electability considerations influence who elites choose to endorse, while competition and uncertainty affect how quickly they coalesce around a single candidate. Some nomination cycles, such as the 2020 Democratic primary, exhibit rapid elite coalescence—where endorsements quickly consolidated behind a frontrunner (e.g., Joe Biden after South Carolina), reinforcing his dominance. Others, such as the 1988 Democratic primary, highlight the consequences of elite fragmentation. That year, Democratic elites failed to coordinate early and split endorsements between multiple candidates, including Michael Dukakis, Jesse Jackson, Al Gore, and Dick Gephardt. The lack of early elite consensus prolonged the nomination battle and prevented the party from consolidating behind a single candidate until later in the cycle, demonstrating how fragmented endorsements can dilute their collective impact.

Despite the strong empirical relationship between endorsements and electoral success, the mechanisms that govern elite coordination and fragmentation remain less well understood. While previous studies have demonstrated that endorsements matter, they have not fully accounted for the strategic dynamics behind elite decision-making. This paper seeks to build on these findings by offering a formal theoretical model that explores the conditions under which elite endorsements act as strategic complements or substitutes, influencing the overall structure of nomination contests.

\paragraph{Theoretical Models of Endorsements}
Early theoretical models primarily treat endorsements as informational signals that help voters update their beliefs about candidate quality. McKelvey and Ordeshook \parencite*{McKelvey195} and Gorfman and Norrander \parencite*{Grofman1990} frame endorsements as tools that provide clarity to voters, assuming that elites act primarily as vehicles of information rather than as strategic players.

Later models take a strategic approach, emphasizing how endorsements shape coordination among voters and elites. Ekmekci \parencite*{Ekmekci2009} models endorsements as a means for uninformed voters to overcome coordination problems but also demonstrates how ideological bias can distort endorsement behavior and allow elites to manipulate voter perceptions. Similarly, Grossman and Helpman \parencite*{Grossman1999} model endorsements as a tool used by interest groups to sway electoral outcomes, rather than as neutral information mechanisms.

Despite these advancements, most existing models assume that endorsements inherently reinforce each other, leading to automatic elite coordination. These models fail to address situations where endorsements might function as strategic substitutes, discouraging further endorsements rather than encouraging them. My work directly fills this gap by formalizing the conditions under which endorsements act as complements or substitutes.

\paragraph{Endorsements and Intraparty Dynamics} Political parties are often viewed as unified organizations that manage internal differences to maximize electoral success. From this perspective, party elites should ideally coordinate their endorsements to signal a clear frontrunner and prevent internal fragmentation. However, while parties have an interest in maintaining electoral viability, they are also composed of individual actors with distinct objectives. Rather than acting as a singular entity, the party becomes a battleground where endorsements reflect power struggles rather than coordinated messaging.

Theoretical models of party organization emphasize that parties serve as mechanisms for regulating factional competition \parencite{Caillaud2002}. Some scholars view parties as informative brands that help voters aggregate signals about ideology and policy direction \parencite{Snyder2002}. Others argue that primary elections themselves are a tool to foster unity by accommodating underrepresented factions. Hortala-Vallve and Mueller \parencite*{Hortala2015} suggest that by relinquishing some control over candidate selection, party elites use primaries as a commitment device to integrate different factions, ensuring their continued loyalty in the general election. From this perspective, endorsements should function as a stabilizing mechanism, reinforcing party unity by guiding voters toward an electable candidate. However, this assumes that elites prioritize long-term party cohesion over short-term strategic positioning.

Endorsements can also function as a mechanism for controlling access to power, rather than simply as a tool for maximizing electoral success. In some cases, party elites strategically coordinate to block outsider candidates and preserve the party’s establishment leadership \parencite{Buisseret2020}. When elites fail to coordinate, outsider candidates can exploit elite divisions, as seen in nomination cycles where the establishment struggles to consolidate support. At the same time, endorsements may reflect intraparty bargaining, where factions exchange political support for influence over policy and candidate selection \parencite{Persico2011}. These dynamics suggest that endorsements are not always reinforcing signals of viability but rather strategic actions aimed at shaping the party’s long-term direction.

The assumption that parties act as unitary actors underlies many theories of party behavior. However, endorsement patterns challenge this assumption by revealing how elite incentives sometimes undermine coordination rather than reinforce it. While Hortala-Vallve and Mueller \parencite*{Hortala2015} emphasize the role of primaries in maintaining unity, their framework assumes that elites remain committed to the party’s collective interest. This paper instead highlights how endorsements emerge from the strategic interactions of individual elites, each weighing ideological, electoral, and personal incentives.
